%!TEX root = ../TM2.tex

% \setlength{\epigraphwidth}{0.40\textwidth}
% \epigraph{Creativity is just connecting things.}{\textit{Steve Jobs}}
 
  \setlength{\epigraphwidth}{0.40\textwidth}
 \epigraph{\large{Eine gute Theorie ist \\ das Praktischste was es gibt.}}{\textit{Gustav Robert Kirchhoff, 1824-1887}}
%\renewcommand\dictumwidth{0.40\textwidth}
%\dictum[\textit{Gustav Robert Kirchhoff, 1824-1887}]{Eine gute Theorie ist \\ das Praktischste was es gibt.}

\vfill
% \hrule
\noindent
In der vorliegenden Version dieser Vorlesungsunterlagen sind noch mehrere Druckfehler enthalten. Für Hinweise an \href{mailto:franz-josef.falkner@mci.edu}{franz-josef.falkner@mci.edu} und/oder \href{mailto:lukas.moschen@gmail.com}{lukas.moschen@gmail.com} sind wir sehr dankbar.
\newpage

%\chapter*{Vorwort}
%\textcolor{red}{FJ - hier werde ich noch ein Vorwort schreiben.}
%
%\emph{Die Kinematik ist die Lehre der Bewegung, ohne deren Ursache zu betrachten. Berücksichtigt man die Trägheit von Körpern, so sprechen wir von der Kinetik}.
%Die beiden Sätze werden gerne in Lehrbüchern verwendet um kurz und knackig die Erkenntnisse zweier Super-Stars der Mechanik zu würdigen.
%
%\prsn{Johannes Kepler} publizierte 1619 sein Werk \emph{Harmonices Mundi} (\emph{Harmonien der Welten}), in welchem Gleichungen zur Beschreibung der Bewegung von Planeten präsentiert wurden - diese kennen wir heute unter den \emph{drei Kepler'schen Gesetzen der Planetenbewegung}.
%Rund 70 Jahre später, im Jahr 1687, veröffentlichte \prsn{Sir Isaak Newton} neben anderen wichtigen Erkenntnissen die \emph{Gravitationstheorie} in seinem Werk \emph{Philosophiae Naturalis Principia Mathematica} (\emph{Mathematische Prinzipien der Naturphilosophie}) und veränderte die Physik grundlegend.
%
%% Warum sind die Erkenntnisse dieser beiden Herren für uns im 21. Jahrhundert immer noch derart wichtig?
%% Aus technischer Sicht ist die sogenannte \emph{klassische Mechanik}, deren Grundstein \prsn{Newton} legte, hinreichend genau um viele physikalische Probleme zu beschreiben.
%Philosophisch betrachtet unterscheidet sich \prsn{Newtons} Gravitationstheorie von allen anderen bisherigen Erkenntnissen der Mechanik.
%\prsn{Kepler} widmete sich der Fragestellung \enquote{\emph{wie} sich Planten bewegen}, darauf aufbauend begründete schließlich \prsn{Newton} \enquote{\emph{warum} sich Planeten bewegen}.
%Erkennen Sie an dieser Stelle die Parallele zu den beiden Säulen der Mechanik, der Kinematik sowie der Kinetik?
%
%Dieses Skriptum ist auf eine ähnliche Weise wie die Abfolge der geschichtlichen Erkenntnisse von \prsn{Kepler} und \prsn{Newton} strukturiert.
%In dem ersten Abschnitt des Skriptums widmen wir uns der Kinematik - wir beschreiben \emph{wie} sich mechanische Systeme bewegen.
%Im zweiten Teil, der Kinetik, befassen wir uns mit der Begründung von Bewegungen, also \emph{warum} Bewegungen stattfinden.
%
%Bitte beachten Sie, dass dieses Skriptum den Besuch der Lehrveranstaltung aus Technischer Mechanik 3 am MCI nicht ersetzt, sondern lediglich ergänzt.
%Ich wünsche Ihnen viel Freude und Erfolgt beim Studium dieses interessanten Fachgebietes!
%\vfill%
%\noindent
%\begin{minipage}[t]{0.40\textwidth}
%	\flushleft Wien, im Februar 2021
%\end{minipage}%
%	\hfill
%\begin{minipage}[t]{0.40\textwidth}
%	\flushright Lukas Moschen
%\end{minipage}%
%\clearpage